\documentclass[9pt,twoside,lineno]{pnas-new}

\templatetype{pnasresearcharticle}

\title{Is correlation meaningfully correlated with causation? Crud, $1/f$ spectra, and an empirical null for small-effect observational claims}

\author[1]{(Author)}
\affil[1]{(Affiliation)}

\leadauthor{(Lead author)}

\significancestatement{Observational science routinely interprets small adjusted correlations as evidence for causal effects. Across nine large datasets, most pairs of measured variables are correlated to a non-trivial degree even after standard adjustments, and the typical small adjusted correlation in observational medicine, nutritional epidemiology, common-variant genetics, and brain-wide association studies has a magnitude not distinctive against the same-domain background. We propose the crud-aware percentile rank, which places any reported correlation within the empirical distribution of same-pipeline residuals. The diagnostic is a magnitude check complementary to identification: a small adjusted correlation whose magnitude is in the bulk of the empirical null is not, by itself, evidence of a direct causal relation, even if its classical $p$-value is tiny.}

\authorcontributions{(To be filled.)}
\authordeclaration{The author declares no competing interests.}
\correspondingauthor{\textsuperscript{1}E-mail: (corresponding author)}

\keywords{causal inference $|$ crud factor $|$ background dependence $|$ $1/f$ spectra $|$ empirical null}

\begin{abstract}
Causal claims in observational research typically interpret an adjusted correlation as evidence of a causal effect. For this logic to work, adjusted correlations must be informative about underlying causal effects. We test this assumption using nine large datasets spanning health, neuroscience, psychology, genomics, politics, and vision. Across all nine, most pairwise correlations are classically significant; the residual crud scale $\sigma_K$ shrinks only slowly under generic adjustment because shared variance is distributed across many components in approximately $1/f$ eigenvalue spectra. Even after removing $K=10$ principal components, $r=0.10$ does not exceed the conventional 95th-percentile magnitude threshold in six of the nine datasets, regardless of sample size: a magnitude unremarkable against same-pipeline background pairs. We propose the crud-aware diagnostic---a percentile rank in the empirical version, a $z$-test in a parametric small-study version---and use it to show that the typical published effect magnitude in observational medicine, nutritional epidemiology, common-variant genetics, and BWAS-style neuroscience does not clear the same-domain background. A decision-theoretic benchmark (Bayes error under a stylized Gaussian mixture) formalizes the magnitude-only limit. The diagnostic is a magnitude check complementary to identification: small adjusted correlations supported by RCT replication or mechanism (e.g., aspirin and cardiovascular mortality) can still encode real causation; magnitudes typical of randomly drawn pairs in the same domain cannot, by themselves.
\end{abstract}

\dates{This manuscript was compiled on \today}
\doi{www.pnas.org/cgi/doi/10.1073/pnas.XXXXXXXXXX}

\usepackage{amsmath,amssymb,amsthm}
\usepackage{booktabs}
\usepackage{longtable}

\newtheorem{theorem}{Theorem}
\newtheorem{lemma}{Lemma}
\newtheorem{assumption}{Assumption}
\newtheorem{corollary}{Corollary}

\newcommand{\E}{\mathbb{E}}
\newcommand{\Var}{\operatorname{Var}}
\newcommand{\Cov}{\operatorname{Cov}}
\newcommand{\SD}{\operatorname{SD}}
\newcommand{\PP}{\mathbb{P}}
\newcommand{\R}{\mathbb{R}}
\newcommand{\diag}{\operatorname{diag}}
\newcommand{\tr}{\operatorname{tr}}
\newcommand{\1}{\mathbf{1}}
\newcommand{\Ph}{\Phi}

\begin{document}

\maketitle
\thispagestyle{firststyle}
\ifthenelse{\boolean{shortarticle}}{\ifthenelse{\boolean{singlecolumn}}{\abscontentformatted}{\abscontent}}{}

\input{introduction}
\input{results_data}
\input{results_theory}
\input{calibration}
\input{discussion}

\matmethods{
\subsection*{Datasets.} Nine large public datasets spanning population health, neuroscience, genomics, psychology, political science, and vision (Section~\ref{sec:datasets-main}; SI Appendix~\ref{si:datasets} gives provenance and dimensions): NHANES, precinct voting + Census, GTEx v8, Allen Institute mouse RNA-seq, HEXACO, CIFAR-10, Kay fMRI, Haxby fMRI, and Stringer mouse V1.

\subsection*{Preprocessing.} Each feature is z-scored (mean zero, unit variance). Sample sizes used in analysis are capped at $n\approx 10{,}000$ per dataset for computational tractability where the raw $n$ is larger.

\subsection*{Adjustment.} The deliberately generic adjustment is regression on the top $K$ principal component scores, with $K=10$ as the headline choice and sensitivity reported across $K\in\{0,1,10,50\}$ (Table~\ref{tab:crud-scale}). PCA is computed via randomized SVD with up to $300$ components. Alternative adjustments (Lasso) are reported in SI Appendix~\ref{app:alt-adjust}.

\subsection*{Empirical null and crud scale.} For each dataset and each $K$, we compute residual correlations across uniformly sampled feature pairs (or, where $\binom{p}{2}\le 5\times 10^5$, all pairs). The crud scale $\sigma_K$ is the SD of these residual correlations. Eigenvalue spectra are fit by power-law on the first 40 components against an exponential alternative via AIC.

\subsection*{Statistical analysis.} The crud-aware $z$-test compares an observed correlation $r$ to $\sigma_K$ via Equation~\eqref{eq:z-crud-main}. The full empirical null replaces the Gaussian formula with the empirical fraction of reference pairs whose absolute adjusted association exceeds the target's. Calibration of the parametric formula against the empirical null in heavy-tailed datasets is reported in SI Appendix~\ref{app:bridge-lemmas}, Table~\ref{tab:heavy-tail-cal}. A reference catalog mapping common application domains to a recommended $\sigma_K$ is in SI Appendix~\ref{si:reference-catalog}.

\subsection*{Code and data.} All analysis is in the Python package at \url{https://github.com/koerding/crud}. An interactive demonstration and a standalone crud-aware test calculator are at \url{https://koerding.github.io/crud/} and \url{https://koerding.github.io/crud/test.html}.

}

\showmatmethods

\acknow{(Acknowledgments to be filled.)}
\showacknow

\bibliography{references}

\end{document}
